% Generated draft: requires mathematical and source review.
\documentclass[11pt]{article}
\usepackage[margin=25mm]{geometry}
\usepackage{fontspec}
\setmainfont{Latin Modern Roman}
\usepackage{amsmath,amssymb,mathtools,unicode-math}
\setmathfont{Latin Modern Math}
\usepackage{enumitem,needspace,xurl,hyperref}
\setlength{\parindent}{0pt}
\setlength{\parskip}{5pt}
\setlength{\emergencystretch}{3em}
\newcommand{\N}{\mathbb{N}}
\newcommand{\Z}{\mathbb{Z}}
\newcommand{\Q}{\mathbb{Q}}
\newcommand{\R}{\mathbb{R}}
\newcommand{\C}{\mathbb{C}}
\newcommand{\F}{\mathbb{F}}
\newcommand{\A}{\mathbb{A}}
\newcommand{\PP}{\mathbb{P}}
\newcommand{\E}{\mathbb{E}}
\newcommand{\eps}{\varepsilon}
\newcommand{\dd}{\,\mathrm{d}}
\begin{document}
\section*{UnsolvedMath ID: 7200071}
\textbf{Source problem code:} catalogue entry
\subsection{Definitions and mathematical statement}
Normalize the common cell volume to one. A geometric-measure formulation considers a locally finite partition $\mathcal E=(E_i)$ of $\mathbb R^3$ into cells of finite perimeter with $|E_i|=1$, pairwise disjoint and covering space up to null sets. Count each internal interface once by
\[
A_R(\mathcal E)=\frac12\sum_i P(E_i;B_R(0)).
\]
One asymptotic area objective is
\[
a(\mathcal E)=\limsup_{R\to\infty}\frac{A_R(\mathcal E)}{|B_R(0)|}.
\]
For periodic partitions, the corresponding density is computed in a period domain and is one-half the average cell surface area per unit cell volume.

Determine the least possible interface area per unit volume among equal-volume space partitions: this is the Kelvin problem. In particular, is the Weaire--Phelan structure globally minimizing? In a fixed admissible-class/density formulation, this asks whether
\[
a(\mathcal E)\geq a(\mathcal E_{\mathrm{WP}})
\]
for every admissible equal-volume partition, where $\mathcal E_{\mathrm{WP}}$ is the precisely specified Weaire--Phelan partition. The source names this structure but gives no construction data. It also does not specify periodicity, the admissible regularity/topology of cells, or the limiting convention for arbitrary nonperiodic partitions, so the displayed density is an explicit possible framework rather than a verified unique source convention.

The supplied summary reports that Weaire--Phelan (1994) found a two-cell-type structure with smaller mean area than Kelvin's truncated-octahedral structure. It describes global Weaire--Phelan optimality as conjectural, with no rigorous complete proof located through its 2026 assessment. No polyhedral-interface restriction or global uniqueness claim is added.
\subsection{Short English statement}
Find the least interface area per unit volume for a partition of three-dimensional space into equal-volume cells, and decide whether the Weaire--Phelan foam attains it. Specify the admissible partitions and asymptotic area objective.
\subsection{Sources}
\begin{itemize}
\item Wikipedia contributors, \emph{List of unsolved problems in mathematics}, Geometry bullet 71, pinned revision 1366636767, \url{https://en.wikipedia.org/w/index.php?title=List_of_unsolved_problems_in_mathematics&oldid=1366636767#Geometry}; accessed 2026-07-29, MediaWiki source extraction. Wikipedia text is CC BY-SA 4.0, with attribution/share-alike handling required for release.
\item Weaire--Phelan (1994), as the two-cell-type improved-foam attribution in the supplied summary; no full publication citation is given. ``On the Kelvin Problem,'' arXiv:1202.1719, is the supplied status/context reference, without authors in the record.
\item Local registry \texttt{data/unsolvedmath/problems.json}, permanent UnsolvedMath ID $7200071$, source code catalogue entry, \emph{The Kelvin problem and optimality of the Weaire--Phelan structure}. Dataset attribution: Ulam AI and the UnsolvedMath Contributors.
\end{itemize}
\end{document}
